\chapter{ Динамическая HP-модель}
\label{ch:chhp}

 

В этой главе рассматривается динамическая HP-модель на квадратной решётке. 
Численное моделирование выполняется в большом каноническом ансамбле методом Монте-Карло, описанным в разделе \ref{MCMC_sokal}. 
В отличие от стандартной HP-модели, здесь меняется не только конформация цепи, но и внутренние состояния мономеров. 
Поэтому энергетический выигрыш зависит от того, какие контакты в данной конформации оказались контактами типа $H$--$H$.

Задача главы --- построить фазовую диаграмму в используемых переменных и качественно локализовать область перехода клубок-глобула.  


\section{Фазовая диаграмма в большом каноническом ансамбле}
 
 
 Для построения фазовой диаграммы динамической HP-модели проводилась серия  вычислительных экспериментов  при различных значениях параметров $\beta$ и $x$, определения которых представлены в разделе описания модели \ref{sec:2_ch:hp}. В качестве основной наблюдаемой величины использовалась средняя длина цепи $\langle N \rangle$.
 
 Для оценки критического значения  химического потенциала при фиксированном $x$ использовался  следующий закон скейлинга:
 \begin{equation}
 	\langle N \rangle \approx \frac{\beta \gamma(x)}{\beta_c(x)-\beta},
 	\label{eq:Nbeta_fit_rus}
 \end{equation}
 который является обобщением выражения для невзаимодействующего блуждания без самопересечений \cite{BerettiSokal}. 
Приближение  численных данных с помощью формулы \eqref{eq:Nbeta_fit_rus} позволяет получить приближённые значения критической линии $\beta_c(x)$.
 
 \begin{figure}%[ht]
 	\centering
 	\includegraphics[scale=0.4]{images/ph_hp.png}
 	\caption{Фазовая диаграмма модели: оценки критических значений $\beta_c$ с погрешностями. Для динамической HP-модели приведены результаты Монте-Карло вычислений, выполненных при $4\times 10^9$ шагах. Типичная длина цепей в этих расчётах составляет от 50 до 200. Для сравнения также показаны значения $\beta_c$ для гомополимера с последовательностью $HH\cdots H$ из работы \cite{Nidras_1996}.}
 	\label{fig:ph_hp}
 \end{figure}
 
 На рисунке \ref{fig:ph_hp} показана полученная фазовая диаграмма. 
 Для динамической модели HP по вертикальной оси откладывается величина $2\beta_c$, что согласуется с предельным случаем отсутствия взаимодействия и позволяет корректно сравнивать результаты с гомополимерной моделью. 
 Для сравнения на тот же рисунок нанесены значения критического химического потенциала  для взаимодействующего блуждания без самопересечений с фиксированной последовательностью $HH\cdots H$ \cite{Nidras_1996}.
 
 Из рисунка видно, что зависимость критического потенциала $\epsilon_a$ от параметра взаимодействия $\epsilon_b$ в двух моделях качественно одинакова. 
 При фиксированном $\epsilon_b$ критический потенциал в динамической HP-модели оказывается больше, чем в гомополимерном случае. 
 Это естественно, поскольку в  изучаемой  динамической HP-модели энергетический выигрыш возникает только для контактов типа $H$-$H$, тогда как пары $P$-$P$ и $P$-$H$ вклада в энергию не дают.
 
 
 
   
 \section{Геометрический переход}
 
 
 Для изучения геометрических свойств цепи анализировалась зависимость среднего квадрата расстояния между концами цепи от её длины \eqref{endtoend}. 
  
 
 
 Так как в численных расчётах доступны лишь конечные длины цепей, для обработки данных использовалась формула для масштабирования \cite{BerettiSokal}: 
 
 \begin{equation}
 	\langle \log(R_N^2+k_1)\rangle = 2\nu \log(N+k_2) + b,
 	\label{eq:nu_fit_rus}
 \end{equation}
 где $k_1$, $k_2$ и $b$ являются подбираемыми  параметрами.
 
 
 
 \begin{figure}%[htb]
 	\centering
 	\includegraphics[scale=0.38]{nu_eps_hp.png}
 	\caption{Оценки метрического показателя $\nu$ в зависимости от $\epsilon_b$. Показаны результаты для динамической HP-модели и для гомополимера. Горизонтальными линиями отмечены значения $\nu=3/4$ для невзаимодействующего  БСП и $\nu=4/7$ для $\theta$-точки. Использованы данные для цепочек длины $N=100$--$300$.}
 	\label{fig:nu_eps_hp}
 \end{figure}
 
 
Результаты оценки показателя $\nu$ представлены на рисунке \ref{fig:nu_eps_hp}. 
По этому графику видно, как меняется геометрия цепи при увеличении $\epsilon_b$. 
При малых $\epsilon_b$ оценки $\nu$ для динамической HP-модели и для гомополимера близки к значению \eqref{nur}
 \[
 \nu=\frac{3}{4},
 \]
 характерному для классического невзаимодействующего блуждания без самопересечений. 
 Это  соответствует развернутой фазе полимерной цепи  --- режим сильного растворителя.
  
 С увеличением $\epsilon_b$ показатель $\nu$ уменьшается, то есть цепь переходит к более компактным конфигурациям.
  На том же рисунке проведена горизонтальная линия, соответствующая значению
  \[
  \nu=\frac{4}{7},
  \]
  которое известно для $\theta$-точки гомополимерной модели \eqref{nu_theta}.
   Для гомополимера пересечение с этим уровнем происходит вблизи   значения $\epsilon_{b\theta}\approx 0.6673(5)$ \cite{Caracciolo_2011}, что подтверждает корректность применяемой процедуры обработки данных.
  
 Если использовать значение $\nu=4/7$ как ориентир для $\theta$-области, то по данным графика \ref{fig:nu_eps_hp} получаем оценку
 \begin{equation}
 	\label{weirdapprox}
 	1.2 \le \epsilon_{b\theta} \le 1.35.
 \end{equation}
 Эта область лежит правее гомополимерной $\theta$-точки. 
 Сдвиг ожидаем: в динамической HP-модели не каждый топологический контакт уменьшает энергию, а только контакт двух гидрофобных мономеров.
  
  \textbf{Оценка положения $\theta$-точки}
  
  В качестве независимого способа определения положения $\theta$-точки была рассмотрена теплоёмкость
  \begin{equation}
  	C=\epsilon_b^2\left(\langle m^2\rangle-\langle m\rangle^2\right),
  	\label{eq:heat_rus}
  \end{equation}
  где $m$ --- число контактов типа $H$-$H$ при фиксированной длине цепи $N$.
  
	  Для конечных длин максимум теплоёмкости является псевдокритическим ориентиром: он сглажен и может смещаться из-за конечномерных эффектов. 
	  Результаты вычисления теплоёмкости для различных $N$ приведены на рисунке \ref{fig:cm}. 
	  С ростом длины цепи максимум становится более выраженным и попадает в ту же область $\epsilon_b$, где на рисунке \ref{fig:nu_eps_hp} показатель $\nu$ приближается к $4/7$.
  
  По положению максимума теплоёмкости получена более узкая оценка для $\theta$-точки:
  \begin{equation}
  	1.25 \le \epsilon_{b\theta} \le 1.27.
  \end{equation}
  Эта оценка согласуется с анализом метрического показателя  \eqref{weirdapprox}.
  
  \begin{figure}[htb]
  	\centering
  	\includegraphics[scale=0.38]{new_heat_c1.png}
  	\caption{Теплоёмкость $C$ как функция $\epsilon_b$ для нескольких длин цепи.   С ростом $N$ максимум теплоёмкости становится более выраженным, что позволяет оценить примерное положение $\theta$-точки.}
  	\label{fig:cm}
  \end{figure}
  
  
  
    \subsection{Вывод}  
    
 
	    В этой главе построена фазовая диаграмма динамической HP-модели в большом каноническом ансамбле и прослежено изменение метрического показателя $\nu$ при росте $\epsilon_b$. 
	    В режиме малых $\epsilon_b$ цепь ведёт себя как невзаимодействующее БСП, а при усилении взаимодействия постепенно переходит к более компактным конфигурациям.
	    
	    Две независимые диагностики дают согласованную локализацию $\theta$-области: по пересечению уровня $\nu=4/7$ получается широкий интервал $1.2\le\epsilon_{b\theta}\le1.35$, а по максимуму теплоёмкости --- более узкая оценка $1.25\le\epsilon_{b\theta}\le1.27$. 
	    По сравнению с гомополимером эта область сдвинута к большим значениям взаимодействия, поскольку энергетически выгодны только контакты $H$--$H$.
	    
	    Численное исследование здесь остаётся качественным. 
	    Ограничения по длинам цепей, связанные с фазовой зависимостью $\beta_c(x)$, и длинные автокорреляционные времена алгоритма Берретти--Сокала не позволяют претендовать на точную реконструкцию критических параметров.
	    Поэтому результаты этой главы используются как мотивация для перехода к каноническому ансамблю, где фиксированная длина цепи позволяет надёжнее анализировать структурные и магнитные переходы в спин-полимерных моделях.
  
  
